\section{Discusión}

\subsection{Respuesta a la pregunta de investigación}

Las estrategias de preprocesamiento modifican la estabilidad de los descriptores
de tres maneras distintas: con un sesgo sistemático (compacidad), con un sesgo que
depende del sitio (relación de aspecto) y con error aleatorio (simetría). Ninguna
estrategia de realce mejoró la línea base de forma consistente entre sitios. Y el
efecto no se limita a los valores: según la estrategia, cambian entre
\FinDiscMin{} y \FinDiscMax{} de las 10 conclusiones entre grupos. El
preprocesamiento puede cambiar las conclusiones científicas que se extraen de los
descriptores, no solo sus valores.

\subsection{Coherencia entre análisis independientes}

\begin{itemize}
    \item \textbf{OE2 y OE3 coinciden.} Las métricas de imagen de OE2 y el Dice de
    OE3 llegan a la misma conclusión por caminos distintos: el realce no mejora la
    línea base. Las imágenes tienen poco ruido, y filtrar o realzar cuesta más de
    lo que aporta.
    \item \textbf{OE2 explica el efecto de sitio.} La firma espectral de cada origen
    se concentra en las subbandas vertical y diagonal de escala fina. Eso explica
    que la energía wavelet difiera entre hospitales, y justifica analizar cada
    origen por separado y no usar el origen como predictor.
    \item \textbf{OE4 anticipa OE5.} El tipo de error predice si una conclusión
    entre grupos sobrevive a la segmentación automática: un sesgo uniforme
    (compacidad) no la altera; un error que difiere entre grupos o sitios
    (relación de aspecto) puede borrar una diferencia real o crear una inexistente.
    \item \textbf{El modelado confirma que el desplazamiento es aprendible.} Los
    descriptores automáticos contienen información suficiente para estimar la
    fracción pulmonar de referencia con un error de unos 3 puntos porcentuales, lo
    que indica un desplazamiento estructurado y no ruido puro.
\end{itemize}

\subsection{Implicaciones}

El criterio para elegir una segmentación debe depender del uso posterior de los
descriptores: la estrategia con mejor solapamiento no fue la que mejor preservó las
comparaciones entre grupos. Toda métrica debe reportarse estratificada por sitio.
Para comparaciones entre sitios, la relación de aspecto y la simetría son poco
confiables; la fracción de área y la compacidad, aunque sesgadas, se comportan de
forma más predecible.

\subsection{Limitaciones}

\begin{enumerate}
    \item \textbf{Convención de trazado.} Las máscaras de Montgomery excluyen el
    corazón y los grandes vasos; las de Shenzhen los delimitan. Es un sesgo
    sistemático entre orígenes que no puede corregirse.
    \item \textbf{Una sola máscara por imagen.} No se puede medir la variabilidad
    entre observadores, y por tanto no se sabe si un error de 0,03 en fracción de
    área es clínicamente relevante.
    \item \textbf{Sin identificador de paciente.} Se agruparon los dos casos
    demostrados por referencia cruzada; 90 imágenes de Montgomery comparten sexo y
    edad con otra sin que eso pruebe identidad.
    \item \textbf{Dos sitios.} No se puede separar el efecto del detector del
    protocolo ni de la población, ni el del detector del de la reducción de
    resolución, porque Shenzhen no documenta el tamaño de píxel. Ningún resultado
    se verificó en un tercer centro.
    \item \textbf{Montgomery pequeño y desbalanceado.} Con pocos pacientes por grupo y
    grupos que difieren en edad y sexo, sus comparaciones son imprecisas y están
    confundidas.
    \item \textbf{Energía wavelet.} Anular la imagen fuera de la máscara crea un
    borde artificial, de modo que el descriptor refleja en parte el contorno.
    \item \textbf{Parámetros declarados.} Los resultados valen para los parámetros
    usados; otros límites de CLAHE u otros cortes de frecuencia no se exploraron.
    \item \textbf{Umbrales convencionales.} El criterio de OE2 (5\,\% y 0,90) y el de
    estabilidad de OE5 son convenciones fijadas antes de ver los datos. Dos
    estrategias quedaron a menos de 0,005 del umbral de bordes, y el criterio de
    estabilidad trata como inestable un cambio de signo entre efectos nulos.
    \item \textbf{Muestra de OE2.} Las comparaciones de OE2 usan 40 imágenes.
    \item \textbf{Etiqueta.} Proviene de lecturas clínicas, no de confirmación
    diagnóstica.
\end{enumerate}

\subsection{Desviaciones respecto al plan}

\begin{itemize}
    \item OE2 estaba previsto para las semanas 6 a 11 y se completó en la 14; el
    cronograma lo registra.
    \item Dos métricas de OE2 se redefinieron (ruido y conservación de bordes), por
    las razones expuestas en la sección \ref{sec:met-oe2}.
    \item La regresión estaba como no planeada; el laboratorio de modelado la exigió
    y se formuló como pregunta complementaria.
    \item El análisis exploratorio declaró el filtrado en frecuencia como no
    pertinente para el flujo; OE2 lo evaluó igual, para cumplir la cobertura
    comprometida y confirmar esa decisión con datos.
    \item OE3 se había evaluado solo en entrenamiento; se completó en validación,
    como pedía la formulación.
\end{itemize}